\chapter{连续时间中的主动推断}
\section{引言}

万物流变，无物恒常不动。

\hfill --- 赫拉克利特，公元前 501 年

本章在第 7 章的基础上继续讨论如何构建生成模型。这里的重点是连续状态空间模型。这类模型尤其适合刻画作用于感觉受体的物理波动，以及我们用来改变周围世界的效应器（例如肌肉）的连续运动。这些模型有许多应用。本章将说明其使用背后的基本原理，突出运动控制中常用的模型类型，以及在此类模型中发挥作用的动力系统，并简要讨论广义同步这一概念。最后，我们将讨论离散生成模型与连续生成模型之间如何得到统一。

\section{运动控制}

正如我们在第 4 章所见，支撑连续时间主动推断的生成模型可以写成一对随机方程，它们决定了状态 $x$ 如何生成数据 $y$，以及状态如何在某个静态变量 $v$ 的作用下随时间演化：

\[
y = g(x) + \omega_y
\tag{8.1}
\]
\[
\dot{x} = f(x,v) + \omega_x
\]

这些方程以及与波动 $\omega$ 相关的精度，共同决定了用于推断感觉原因的模型。

注意，动作并未出现在公式 8.1 中。这是因为，如第 6 章所述，动作属于生成过程，而不属于生成模型。生成模型只处理那些会被马尔可夫毯之外的状态直接影响的变量。如果我们要写出真实世界的动力学，也就是生成过程，就必须把动作 $u$ 包含进来：

\[
y = g(x) + \omega_y
\tag{8.2}
\]
\[
\dot{x} = f(x,u) + \omega_x
\]

需要注意的是，用来定义生成模型（公式 8.1）的函数 $g$ 与 $f$，以及 $\omega$ 的精度，并不一定与定义生成过程（公式 8.2）时使用的函数相同。正如我们在第 2--4 章中看到的，动作会改变感觉数据，使自由能最小化。这意味着，我们不需要在生成模型中显式写出动作的动力学，它们会从公式 8.1 中各项的具体选择里自然涌现。为了建立直觉，我们从一种非常简单的生成模型开始：

\[
g(x) = x
\tag{8.3}
\]
\[
f(x,v) = v - x
\]

公式 8.3 表明，隐藏状态表示数据的期望值，并且它服从一个简单的，也就是点吸引子的动力学。所谓吸引子，是指当 $x < v$ 时，$x$ 的期望变化率为正；反之则为负。这意味着 $x$ 总会流向 $v$，也就是 $v$ 是一个吸引点或固定点。为了生成数据，我们定义一个简单的生成过程：

\[
g(x) = x
\tag{8.4}
\]
\[
f(x,u) = u
\]

在最小化自由能时，这意味着动作会发生变化，以满足公式 8.3 的预测。如果 $\mu$ 是 $x$ 的期望值，那么能够最小化预测数据 $g(\mu)$ 与观测数据 $y$ 之间差异的动作，就是令 $u = v - \mu$。这正是“平衡点假说”的一种表达（Feldman and Levin 2009）。该假说认为，运动控制可由反射弧实现，肢体只需被简单地拉向由下行运动信号设定的平衡点。在主动推断框架下，这些信号是预测，更具体地说，是本体感觉预测，例如对肢体或眼睛预期位置的预测（Adams, Shipp, and Friston 2013）。因此，运动控制就是通过动作实现对这些（本体感觉）预测的满足，如图 8.1 的示意所示。值得注意的是，这一方案并不需要规定“逆模型”，也就是那种从期望后果映射到实现该后果所需运动指令的模型；而这类逆模型在其他运动控制理论中被广泛使用（Wolpert and Kawato 1998）。

公式 8.3 所表达的是我们可以在生成模型中使用的最简单吸引子系统。但在许多情形下，它过于简单，因为更现实的情况需要服从牛顿动力学。更复杂的模型会认识到，由肌肉产生的力改变的是速度，也就是引起加速度，而不是直接改变位置。公式 8.5 明确说明了这一点，其中 $x_1$ 是位置，$x_2$ 是速度：

\[
f(x,v) =
\begin{bmatrix}
x_2 \\
\frac{\kappa}{m}(v - x_1)
\end{bmatrix}
\tag{8.5}
\]

这个表达式等价于服从胡克定律的弹簧动力学。位置变化率（第一项）就是速度；速度变化率（第二项）与当前位置和点 $v$ 之间的距离成正比，其比例常数由物体质量 $m$ 与弹簧常数 $\kappa$ 的比值给出。若两边同乘质量，则可得到由连接点 $v$ 与 $x_1$ 的弹簧所产生的力\footnote{原文在此处标注脚注编号 1。}，即 $\kappa(v - x_1)$，它等于质量乘以速度变化率。这正是牛顿第二定律。换言之，我们可以写下一个生成模型，预测如果有一根弹簧把肢体拉向某个期望位置，会出现怎样的动力学。通过预测这种牛顿力学所导致的（本体感觉）数据，我们就能够实施满足这些预测的运动。

\section{动力系统}

如 8.2 节所述，主动推断的连续时间形式特别适合刻画运动。更一般地说，它适合用于指定非线性动力系统的生成模型，在这类系统中，对时间和空间进行离散化会显得低效。最简单的动力系统形式就是公式 8.3 中的吸引子，但更复杂的系统可以产生丰富得多的行为。受本书篇幅限制，我们无法充分介绍基于更复杂动力系统的大量模型研究工作，不过可参见表 8.1 中列出的一些关键进展。这里我们聚焦于理解这些系统所需的若干基本原理。本节将简要概述两类被用于构造此类生成模型的动力系统：Lotka--Volterra 动力学和 Lorenz 系统。前者可以用于描述那些动力学具有顺序性的系统，后者则代表混沌系统。

Lotka--Volterra 动力学起源于生态学中对捕食者--猎物关系的刻画。虽然它后来被应用到许多学科，但捕食者--猎物系统仍是帮助我们形成直觉的有用例子。当捕食者数量较少时，猎物数量可以增长到相对较大的规模。这又为捕食者提供了更多食物，从而使捕食者数量增长。捕食增加会使猎物数量下降，进而又导致捕食者数量下降。随后，这一循环继续进行。于是就会出现一种振荡模式：先是猎物种群达到峰值，然后是捕食者种群达到峰值，接着又轮到猎物，如此往复。如果把这一机制推广到两个以上的种群，例如食肉动物、食草动物和植物三个种群，我们就可以生成一串峰值序列。图 8.2 展示了三个种群下的广义 Lotka--Volterra 动力学，其形式如下：

\[
f(x,v) = x \circ (v + Ax)
\tag{8.6}
\]

这里的 $x$ 与前面一样是一个向量；符号 $\circ$ 表示按元素相乘。向量 $v$ 给出内禀的出生率与死亡率，矩阵 $A$ 的元素则表示不同物种之间的关系：若矩阵中某列索引的物种捕食某行索引的物种，则对应元素为正；若关系相反，则为负。

图 8.2 清楚表明，只要生成模型中包含 Lotka--Volterra 动力学，就能够实现时间顺序结构（Huerta and Rabinovich 2004），具体取决于当前哪个峰值最高。每一条曲线都可以被看作一个隐藏状态，而不必真的对应一个物种。图 8.3 突出了两种利用这类动力学来生成行为的重要例子。其一是一个层级模型，它利用 Lotka--Volterra 系统所提供的顺序动力学，来安排条件刺激相对于眨眼反应的时序（Friston and Herreros 2016）。这一范式基于用于研究小脑功能的实验：一个无条件刺激，例如朝动物眼睛吹气，会引发眨眼反应；如果多次在无条件刺激之前播放一个条件刺激，例如听觉音调，那么动物就可以通过学习（见框 8.2）Lotka--Volterra 动力学中条件刺激与无条件刺激之间相隔的峰值个数，学会预先眨眼，以恰当时机应对气流刺激。这是一种时间学习，因为峰值个数为条件刺激到无条件刺激之间的时间间隔提供了一个隐式估计。

图 8.3 的第二个例子中，每个顺序峰值都与一个不同的吸引点相关联，这些吸引点驱动动作依次移动到一系列位置，而这些位置被安排成类似手写笔画的轨迹（Friston, Mattout, and Kilner 2011）。这两个例子表明，广义 Lotka--Volterra 系统能够利用连续动力系统来提供对顺序动力学的有效建模。

\subsection{框 8.1：精度、注意与感觉衰减}

我们在第 7 章中已经讨论过精度的重要性，但仍有必要回顾它在连续时间系统中的作用。在许多方面，这一概念在连续时间背景下更自然，因为变量 $\Pi$ 直接作为拉普拉斯近似的结果出现。它在推断动力学中直接充当乘性增益（见图 8.1），不同的精度会为信念更新中的不同影响因素赋予不同权重。

将精度解释为突触增益，使其与若干重要的神经生物学问题联系起来。从经验角度看，更高的精度意味着更强烈的信念更新。在电生理研究中，这可能表现为振幅更大且峰值更早的诱发反应；在单细胞记录中，则可能表现为当刺激落入某个神经元感受野时，该刺激对神经元放电率产生乘性影响。这些现象常常与注意加工联系在一起，即某一个感觉通道，或某些通道的子集，相对于其他通道被优先处理。从主动推断的角度看，精度与注意是同义词。利用这一点，人们已经在计算机仿真中复现了多种注意现象，包括 Posner 范式（Feldman and Friston 2010）。具体而言，如果用一个线索来预测来自两个位置之一的感觉输入精度，那么模型就能够复现实证发现：位于被提示位置的刺激，其反应速度快于位于另一位置的刺激。

精度控制的第二个重要方面，是它在运动生成中的作用。为了理解这一点，不妨设想缺乏这种控制会发生什么。首先设想感觉数据以高精度被预测。这样一来，来自这些数据的信息就被赋予较高的突触增益，从而产生关于身体某部分位置的真实推断。问题在于，在主动推断中，运动指令与预测是等价的。一个准确的“我没有在动”的信念，无法用来预测运动的感觉后果，而这种预测恰恰是启动运动所必需的。在高精度感觉输入下，“我正在移动”这一信念会立刻在相反证据面前被纠正，因此动作无法执行。这告诉我们一个重要事实：为了生成运动，我们必须能够暂时忽略该运动的感觉后果，从而形成一个起初并不为真的信念，即“我正在移动”。一旦这个信念建立起来，运动的本体感觉及其他感觉后果就可以被预测，并通过图 8.1 所示的机制付诸实现。对相反证据的这种忽略被称为“感觉衰减”，它代表了运动得以发生所需的精度下降（Brown, Adams et al. 2013; Pezzulo 2013; Seth 2013; Pezzulo, Rigoli, and Friston 2015; Seth and Friston 2016; Allen et al. 2019）。显然，在两次运动之间恢复这种精度是有益的，因为这样才能基于感觉输入进行适当推断。这意味着一个周期性过程：先衰减，再运动；例如扫视期间视觉输入被周期性抑制，随后扫视本身又被抑制。把运动建立在注意暂时中止之上的观点，与一种起源于 19 世纪的意念运动理论密切相关，该理论最初用于解释催眠状态下诱发的动作。

\subsection{框 8.2：连续模型中的学习}

如第 7 章所述，学习是对生成模型参数 $\theta$ 的信念进行优化的过程。在连续时间域中，这意味着要随时间积累证据。其工作方式可以理解为：把一系列无穷小时间区间中的数据视作满足 i.i.d.（独立同分布）假设，并构造一个由时间不变参数生成观测的生成模型：

\[
\ln p(y,\theta) = \ln p(\theta) + \int \ln p(y(t)\mid \theta)\, dt
\]
\[
\approx \ln p(\theta) - \int F[y(t)\mid \theta]\, dt
\]

这可以用来构造一个泛函 $S$，它通过对条件于参数的自由能做时间积分，扮演参数自由能的角色。采用拉普拉斯近似，可得

\[
S(\theta) = \mathbb{E}_{q(\theta)}\!\left[\ln q(\theta) + \int F[y(t)\mid \theta]\, dt - \ln p(\theta)\right]
\]
\[
\approx \int F[y(t)\mid \mu_\theta]\, dt - \ln p(\mu_\theta)
\]
\[
\dot{\mu}_\theta = \partial_{\mu_\theta} S(\mu_\theta)
\]
\[
= \partial_{\mu_\theta}\ln p(\mu_\theta) - \partial_{\mu_\theta}\int F[y(t)\mid \mu_\theta]\, dt
\]
\[
= \partial_{\mu_\theta}\ln p(\mu_\theta) - \alpha
\]
\[
\alpha = \partial_{\mu_\theta} F[y(t)\mid \mu_\theta]
\]

其中，$\alpha$ 起到积累自由能梯度，也就是证据梯度的作用。

第 7 章中的 POMDP 形式，在主动推断应用中很大程度上已经取代了广义 Lotka--Volterra 系统的使用。不过，把这种动力学记作一种可能支撑第 7 章离散顺序动力学的连续系统，仍然是有益的。此外，Lotka--Volterra 系统还明确区分了两类表征：一类是 temporally deep planning 中与序列有关的表征；另一类则是在广义运动坐标中对变化率的表征（见第 4 章）。二者各有位置，但处理的是不同类型的问题。

第二类在主动推断研究中被广泛应用的动力系统是 Lorenz 系统：

\[
\dot{x} =
\begin{bmatrix}
\sigma(x_2 - x_1) \\
x_1(\rho - x_3) - x_2 \\
x_1x_2 - \beta x_3
\end{bmatrix}
\tag{8.7}
\]

这些参数分别被称为 Prandtl 数 $\sigma$、Rayleigh 数 $\rho$，以及一个与系统物理性质相关的常数 $\beta$。根据这些参数取值的不同，系统可能表现出截然不同的行为。Lorenz 吸引子最初是为了解释大气对流动力学而提出的。其漫游式行为促使人们将其纳入生成模型，用于模拟具有挑战性的推断问题。一个重要例子是鸟鸣模拟，下一节将详细展开。此类系统也被用来模拟简单物理系统，并研究在何种条件下其行为开始表现出类似有感知主体的特征。图 8.4 展示了 Lorenz 系统在一组示例参数下的行为。

\section{广义同步}

如上所述，连续状态空间模型的一个关键应用例子，是一系列基于合成鸟鸣的研究（Friston and Frith 2015b）。这些研究的一个重要方面是通信与多主体推断问题。其核心思想在于：一个生物体能够使其内部状态与外部世界中的某物同步，这本质上就是推断。当外部之物是另一个拥有相似模型的生物体时，这种同步就意味着，一个生物体的内部状态应当逐渐变得与另一个生物体的内部状态相似。这可以被视为一种原始形式的心智理论。

图 8.5 展示了用于模拟鸣禽的生成模型。在这个层级模型中，高层状态（第 2 层）按照一个缓慢的 Lorenz 系统演化。该系统的一个维度又被用来参数化低层，也就是第 1 层中一个更快 Lorenz 系统的 Rayleigh 数。随后，低层变量映射到感觉性的声谱图数据。与图 8.1 类似，生成过程还额外包含动作；只不过这里动作并不是移动肢体，而是影响鸣禽的发声器官，从而使鸟能够影响声谱数据。与前面一样，动作被生成出来以消除预测误差。这意味着，如果一只鸟听到了它正在预测的歌曲，它就没有必要自己发声；但如果它预测某段歌声却没有听到，那么它就必须开始歌唱，以消除误差。

当场景中有两只鸟，而且它们的生成模型具有相似结构时，这种动力学就会变得更有趣。只要一只鸟在唱，另一只就不需要唱，因为没有需要消除的误差。然而，如果一只鸟停止歌唱，另一只就必须把同一首歌继续唱下去。这就产生了一种轮流接唱的形式，有时被表述为“唱同一本圣歌本上的歌”，即每只鸟都贡献同一首歌中的一部分。是什么导致了这种轮替？为什么一只鸟不会把整首歌都唱给同类听？答案与感觉衰减有关（见框 8.1）：发声以产生鸟鸣需要降低对动作后果预测的精度。正如在扫视眼动中一样，这意味着主体必须在两种状态之间交替：一是注意感觉数据（视觉或听觉数据），二是在执行旨在改变这些数据的动作期间对感觉影响进行衰减。当系统中涉及两个主体时，这就导致了在倾听对方与自己歌唱之间的交替，也就是一种简单形式的对话。

要使这种合成对话顺利工作，关键在于两只鸟必须彼此同步，并知道自己当前处于歌曲，也就是对话轨迹的哪个位置。这意味着，两只鸟在生成模型中对隐藏状态的推断必须保持一致。图 8.5 右上角展示了一个同步流形，其对象是两只尚未针对彼此优化自身生成模型的鸟；该图描绘了每只鸟对高层隐藏状态所持信念的轨迹。所谓同步，意味着当一只鸟推断出某个特定隐藏状态取值时，另一只鸟也应推断出相同的值；因此，我们会预期这条轨迹贴近 $x = y$ 这条直线，从技术上说这叫作混沌系统的相同同步。围绕这条直线的波动意味着同步不完美，图中的确如此。然而，在两只鸟彼此接触，并学习了对方生成模型的参数之后（见框 8.2），同步就几乎完美了，如图右下所示。这意味着，每只鸟都已经学会了关于另一只鸟的内容，并且能够推断对方“脑中正在发生什么”。简而言之，它们学会了共享同一个叙事，也就是“唱同一本圣歌本上的歌”。

更一般的同步形式并不要求行为必须沿着 $x = y$ 这条线同步。在广义同步中，联合行为所占据的空间维度更低，是一个低维空间，而不是原本可能占据的更高维空间。具体到这里，它占据的是低维的一维空间，而不是更高的二维空间。不过，这个低维空间，也就是同步流形，可以是弯曲的，也可以具有其他形状。这类似于我们虽然生活在二维的地表上，但这个地表本身其实弯曲成一个三维球面。除了在社会行为中的核心作用之外，广义同步，即在高维联合空间中占据一个低维区域，对于把生物系统描述为在进行推断也极其重要；此时所谓推断，正是内部状态与外部状态之间的广义同步。尽管本书没有篇幅深入展开这一庞大主题，但从推断视角出发，我们可以理解某些神经精神综合征，例如自闭症，与广义同步失效之间的关系。需要指出的是，这种同步不仅对连续时间模型重要，在涉及多个主体语言交流的 POMDP 模型中同样重要（Friston, Parr et al. 2020）。

\section{混合模型（离散与连续）}

正如本章以及上一章所展示的，离散模型和连续模型在主动推断中都具有重要应用。虽然许多场景只需要其中一种，但更整体的观点会承认，两者很可能是同时发挥作用的。这就要求我们找到一种方式来结合这两类生成模型，使得单一模型同时包含连续变量与离散变量（Friston, Parr, and de Vries 2017）。这种混合模型能够一方面对顺序性的动作计划进行推断，另一方面又通过连续模型把这些决策转化为具体动作。图 8.6 展示了这类模型的形式：其高层是一个 POMDP，行为方式与第 7 章中描述的一样；而在每一个离散时间步上，它都会生成一个本章所讨论类型的连续模型。这样一来，连续时间就被分解成一串离散排列的短连续轨迹。

为了把离散层的结果翻译到连续层，我们需要把每一个备选结果与某个连续空间中的一点联系起来。为了帮助建立直觉，我们考虑一个延迟期眼动任务的例子，这种任务在灵长类研究中很常见，如图 8.7 所示。该任务包含三个阶段。首先，在猴子保持注视中央注视十字的同时，一个目标会出现在例如四个可能位置中的某一个。接着，目标消失，猴子必须在一段延迟期内记住它的位置。最后，系统给出一个信号，提示猴子执行一次扫视，此时猴子必须把目光移向原始目标出现的位置。要完成这一任务，猴子必须能够对序列进行推断，例如当前处于任务的哪个阶段；同时还要推断四个位置中哪一个才是应该瞄准的位置。这些都是典型的范畴式推断问题，适合用 POMDP 来表述。然而，一旦正确位置被选定，猴子还必须执行一段眼球运动，把中央凹移动到目标位置的连续坐标上。

图 8.7 展示了一个解决该问题的混合生成模型是如何工作的。在顶部图中，模型高层做出范畴性决策：它计算在四个离散时间段上，对四个离散目标位置的后验信念。在中部和底部图中，模型低层则根据高层的离散推断，计算连续的行为轨迹，也就是眼动轨迹。

把关于离散目标位置的决策转换成连续眼动，需要使每一个离散目标位置 $o$ 都对应于一个关于连续隐藏原因 $v$ 的分布，而该分布给出目标的坐标。随后，关于目标坐标的先验可以通过对这些位置进行贝叶斯模型平均而得到，并以 POMDP 层中的推断结果为权重：

\[
P(v \mid o_\tau) = \mathcal{N}(\eta_v^{\,o_\tau}, \Pi_v)
\]
\[
P(v) \approx \mathcal{N}(\eta_v,\Pi_v)
\tag{8.8}
\]
\[
\eta_v = \mathbb{E}_{Q(o_\tau)}[\eta_v] = \sum_i o_{\tau i}\,\eta_v^{\,i}
\]

为了推断哪个离散目标最能解释连续数据，我们还必须能够计算与每一个假设目标对应的证据，而该证据又是观测到的连续数据的函数。这体现出一个事实：混合模型要求层级中的高层与低层之间存在双向相互作用。正如图 8.7 所示，这种结构使系统既能形成“在何处执行扫视”的信念，又能在恰当时刻把这些扫视真正执行出来。在第一个离散时间步，也就是连续时间尺度上的前 250 ms 内，猴子能够较有信心地推断：在第一个时间步上，它的眼睛位于中央注视十字；在第二个时间步上，也就是 500 ms 之前，它将继续保持这种注视；而在这之后，最可能的行动过程将使其中央凹对准上方目标位置。这一点可以从顶部图中的离散放电率看出来。随后，这些离散推断会转化为连续行为；而当这些行为真正被实施时，又会反过来增强关于离散状态的信念。注意，在第三个时间步，也就是 500 到 750 ms 之间，一旦连续数据出现，向上扫视这一选项的概率就会增大。这个基于实验认知研究的简单例子，说明了在离散动作计划与其连续实施之间进行转换的基本原则。

\subsection{框 8.3：混合模型与聚类}

在主动推断之外，把范畴性生成模型与连续生成模型结合起来的问题，主要是通过聚类的视角来理解的。在这种问题中，目标是把每个连续数据点分配到一个离散簇。为解决这一问题，人们已经采用了多种算法，但其中大多数在隐含层面上都依赖于与这里类似的生成模型，也就是高斯混合模型：

\[
P(y,s,D,\eta,\Pi) = P(D)P(\eta)P(\Pi)\prod_i P(s_i \mid D)\,P(y_i \mid s_i,\eta,\Pi)
\]
\[
P(D) = \mathrm{Dir}(d)
\]
\[
P(s_i \mid D) = \mathrm{Cat}(D)
\]
\[
P(y_i \mid s_i,\eta,\Pi) = \mathcal{N}(\eta_{s_i}, \Pi_{s_i})
\]

在聚类方法中，问题是推断每个簇的均值与精度，分别记为 $\eta$ 与 $\Pi$，以及每个数据点 $y_i$ 属于给定簇的后验概率 $P(s_i \mid y_i)$。对于我们当前的目的，也就是 8.5 节中描述的场景，我们假定对 $\eta$ 和 $\Pi$ 具有精确的 delta 函数先验，并通过贝叶斯模型约减来计算 $Q(s_i) \approx P(s_i \mid y_i)$（见框 7.3）。

\section{小结}

本章概述了主动推断框架下连续时间生成模型的应用。这是一个极其庞大的主题，许多内容都未能展开，进一步阅读可参见表 8.1。不过，这里阐述的总体概念已经为继续深入这些模型提供了基础。具体来说，我们从“满足预测”的角度讨论了运动生成。这种处理方式极大简化了运动控制问题，因为我们不再需要额外的机制或逆模型，只需要脊髓或脑干反射弧即可。我们还强调了精度与感觉衰减在这类运动控制中的作用。

鉴于连续方案的一项关键优势，在于能够用动力系统来表述生成模型，我们概述了两类在主动推断研究中被广泛使用的动力系统。广义 Lotka--Volterra 系统为连续背景下的时间排序提供了机制，而 Lorenz 吸引子则能够用来生成丰富的模拟，包括合成鸟鸣。随后，我们讨论了广义同步这一概念。一个系统的内部状态与外部状态之间的同步，构成了以推断方式理解脑功能的基础；同时，它对于社会系统的理论也至关重要，因为在社会情境中，外部状态在很大程度上由同类个体，也就是“像我这样的生物”构成。

最后，我们说明了如何把第 7 章的离散模型与第 8 章的连续模型统一起来，把 POMDP 形式下最小化期望自由能的动力学，包括利用性与探索性行为，与通过连续过程来实现这些行为，以及在两者之间进行调节的双向消息传递机制结合起来。简言之，这一整合让我们能够从决策走向动作，再从动作返回决策。
